% Siegel 定理证明版
\documentclass[11pt,a4paper]{article}

\usepackage{ctex}
\usepackage{amsmath,amssymb,amsthm,mathtools}
\usepackage[margin=1in]{geometry}
\usepackage{hyperref}
\usepackage{enumitem}

\newtheorem{theorem}{定理}
\newtheorem{lemma}{引理}
\newtheorem{proposition}{命题}
\newtheorem{remark}{注}

\newcommand{\OK}{\mathcal O_K}
\newcommand{\Q}{\mathbb Q}
\newcommand{\PP}{\mathbb P}
\newcommand{\Spec}{\operatorname{Spec}}
\newcommand{\Norm}{N_{K/\Q}}

\title{IMO 1988 第 6 题的推广的证明}
\date{}
\begin{document}
\maketitle

\section{命题陈述}

设 \(K\) 是全实数域，\(\OK\) 为其整数环。
固定一个实嵌入 \(K\hookrightarrow\mathbb R\)，称为恒等嵌入。
对 \(x\in\OK\)，若其在该嵌入下的像为正，则记 \(x>0\)。

考虑性质
\[
(P_K):\qquad
\forall a,b\in\OK,\quad
a,b>0,\quad ab+1\mid a^2+b^2
\Longrightarrow
\frac{a^2+b^2}{ab+1}\in \OK^2.
\]

我们证明：
\[
(P_K)\quad\Longleftrightarrow\quad K=\Q.
\]

蕴含 \(K=\Q\Rightarrow(P_K)\) 正是经典的 IMO 1988 第 6 题。
新的方向用 Siegel 定理证明。

\section{关键的单位构造}

设 \(u\in\OK^\times\) 是一个单位。令
\[
a=u^2-1,\qquad b=u^2+1.
\]
则
\begin{align*}
ab+1
&=(u^2-1)(u^2+1)+1\\
&=u^4,
\end{align*}
因此
\[
ab+1\mid a^2+b^2.
\]
更精确地，
\begin{align*}
a^2+b^2
&=(u^2-1)^2+(u^2+1)^2\\
&=2(u^4+1).
\end{align*}
因此 \((P_K)\) 中的商为
\[
q_u
=
\frac{a^2+b^2}{ab+1}
=
2(1+u^{-4}).
\]

由于 \(u^4=(u^2)^2\) 本身在 \(\OK\) 中就是平方，我们有
\[
q_u\in\OK^2
\quad\Longleftrightarrow\quad
2(u^4+1)\in\OK^2.
\]
确实，
\[
q_u=z^2
\Longrightarrow
2(u^4+1)=u^4z^2=(u^2z)^2,
\]
反之
\[
2(u^4+1)=w^2
\Longrightarrow
q_u=(wu^{-2})^2.
\]

因此，不能产生反例的单位恰好就是仿射曲线
\[
C:\qquad w^2=2(u^4+1)
\]
上的整点。

\section{相关的仿射曲线}

考虑
\[
f(u)=2(u^4+1)\in K[u].
\]
在 \(\overline K\) 上，该多项式有四个互异根，因为
\[
f'(u)=8u^3
\]
而 \(u=0\) 不是 \(f\) 的根。
因此仿射曲线
\[
C:\quad w^2=f(u)
\]
是非奇异的。

其光滑射影完备化 \(\widetilde C\) 是一条四次超椭圆曲线。
对于无平方因子的 \(2g+2\) 次多项式，相应的光滑射影曲线亏格为 \(g\)。因此
\[
\deg f=4
\qquad\Longrightarrow\qquad
g(\widetilde C)=1.
\]

因此 \(C\) 是一条正亏格的仿射曲线。

由 Siegel 定理， \(\widetilde C\) 上在 \(\OK\) 上的整点有限

即只有有限多个 \(u\in\OK\) 使得
\[
2(u^4+1)
\]
是 \(\OK\) 中的平方。

定义
\[
E=
\left\{
u\in\OK^\times:
2(u^4+1)\in\OK^2
\right\}.
\]
则
\[
E\text{ 是有限集}.
\]

\section{无穷多个合适的单位}

现在假设
\[
K\neq\Q.
\]
由于 \(K\) 是全实数域，
\[
[K:\Q]=n>1.
\]
Dirichlet 单位定理给出
\[
\operatorname{rank}\OK^\times=r_1+r_2-1=n-1>0.
\]
所以 \(\mathbb Z\) 是 \(\OK^\times\) 的子群，因此 \(\OK^\times\) 是无限群。

取一个无穷阶单位 \(v\)。若 \(v\) 在有些实嵌入下不为正，由嵌入的性质，用 \(v^2\) 代替 \(v\) 即可在所有实嵌入下为正。
由于 \(v\) 是无穷阶单位，所以不存在正整数 \(m\)，使得 \(v^m=1\)，所以在恒等嵌入下的绝对值不被迫为 \(1\)。
若 \(|v|>1\)，则幂 \(v^m\) 给出任意大的正单位。若 \(|v|<1\)，则改用 \(v^{-1}\)。

因此有无穷多个单位
\[
u\in\OK^\times
\]
满足
\[
u>1.
\]

由于 \(E\) 有限，我们可以选取
\[
u\in\OK^\times,\qquad
u>1,\qquad
u\notin E.
\]

\section{反例的构造}

对这个单位 \(u\)，定义
\[
a=u^2-1,\qquad b=u^2+1.
\]
因为 \(u>1\)，
\[
a>0,\qquad b>0.
\]
此外，
\[
ab+1=u^4,
\]
所以当然有
\[
ab+1\mid a^2+b^2.
\]

相应的商为
\[
\frac{a^2+b^2}{ab+1}
=
2(1+u^{-4}).
\]
若这个商在 \(\OK\) 中是平方，则如上所示，
\[
2(u^4+1)
\]
将是 \(\OK\) 中的平方。这将蕴含
\[
u\in E,
\]
与 \(u\) 的选取矛盾。

因此
\[
\frac{a^2+b^2}{ab+1}\notin\OK^2.
\]
于是 \((P_K)\) 不成立。

我们证明了
\[
K\neq\Q
\quad\Longrightarrow\quad
\neg(P_K).
\]

\section{有理数情形}

当 \(K=\Q\) 时，\(\OK=\mathbb Z\)。此时 \((P_K)\) 恰好就是经典的 IMO 1988 第 6 题：
\[
a,b\in\mathbb Z_{>0},\quad ab+1\mid a^2+b^2
\quad\Longrightarrow\quad
\frac{a^2+b^2}{ab+1}\in\OK^2.
\]
下面给出韦达跳跃证明。

设 \(a,b\in\mathbb Z_{>0}\) 满足 \(ab+1\mid a^2+b^2\)，令
\[
k=\frac{a^2+b^2}{ab+1}\in\mathbb Z_{>0}.
\]
则
\[
a^2+b^2=k(ab+1).
\]
假设 \(k\) 不是完全平方。在所有满足上式且 \(k\) 非平方的正整数对 \((a,b)\) 中，取 \(a+b\) 最小的一对，并不妨设 \(a\ge b\)。

将上式看作关于 \(x\) 的二次方程
\[
x^2-kbx+(b^2-k)=0.
\]
它有一个根 \(x=a\)。设另一根为 \(a'\)。由韦达定理，
\[
a+a'=kb,\qquad aa'=b^2-k.
\]
因此
\[
a'=kb-a\in\mathbb Z,\qquad a'=\frac{b^2-k}{a}.
\]

若 \(a'<0\)，则 \(a'\le -1\)，从而
\[
b^2-k=aa'\le -a,
\]
即 \(k\ge a+b^2\)。另一方面，由 \(a\ge b\ge1\)，
\[
k=\frac{a^2+b^2}{ab+1}
<\frac{a^2+b^2}{ab}
=\frac ab+\frac ba
\le \frac ab+1
\le a+1.
\]
于是 \(k\le a\)，这与 \(k\ge a+b^2\) 矛盾。故 \(a'\ge0\)。

又
\[
a'=kb-a<\left(\frac ab+1\right)b-a=b,
\]
而 \(a\ge b\)，所以 \(a'<b\le a\)。若 \(a'=0\)，则
\[
aa'=b^2-k=0,
\]
故 \(k=b^2\) 为完全平方，矛盾；因此 \(a'>0\)。

于是 \((a',b)\) 也是满足条件的正整数对，且对应商仍为 \(k\)，但
\[
a'+b<a+b,
\]
这与 \((a,b)\) 的最小性矛盾。因此 \(k\) 必为完全平方数。
因此
\[
(P_{\Q})\text{ 成立}.
\]

合并两个方向得到
\[
\OK\text{ 具有 }(P_K)
\iff
K=\Q.
\]

\end{document}